\documentclass[11pt]{article}
\usepackage[margin=25mm]{geometry}
\usepackage{fontspec}
\setmainfont{Latin Modern Roman}
\usepackage{amsmath,amssymb,mathtools}
\usepackage{unicode-math}
\setmathfont{Latin Modern Math}
\usepackage{xurl,hyperref}
\hypersetup{colorlinks=true,urlcolor=blue}
\setcounter{secnumdepth}{0}
\setlength{\parindent}{0pt}
\setlength{\parskip}{5pt}
\setlength{\emergencystretch}{3em}
\newcommand{\sref}[2]{\hyperlink{source:#1:#2}{[S#2]}}
\newcommand{\cataloguescope}[1]{\paragraph{Recorded target.}#1}
\newcommand{\R}{\mathbb R}
\newcommand{\C}{\mathbb C}
\newcommand{\N}{\mathbb N}
\newcommand{\Z}{\mathbb Z}
\newcommand{\Q}{\mathbb Q}
\newcommand{\F}{\mathbb F}
\newcommand{\eps}{\varepsilon}
\begin{document}
\section{Physical Regularity Criteria for Massless Scattering Theory}\label{statement}

\subsection{Definitions and mathematical statement}

\paragraph{Shared notation.}$\mathbb N=\{1,2,\ldots\}$; the standard number fields have their usual meanings. All additional conventions are stated locally.


\paragraph{Definitions and conventions.} Haag--Ruelle scattering theory constructs asymptotic multiparticle states from local observables in a relativistic quantum field theory with a vacuum and translation spectrum. A sharp-mass particle has a one-particle mass shell even when other excitations accumulate near it. Herbst's regularity assumption controls the unwanted spectral component near this mass shell sufficiently strongly for the scattering limits. The report concerns a theory that also has massless particles, while the scattered massive particles retain sharp masses.

\paragraph{Source scope and qualifications.} The report excludes electrically charged infraparticles through the sharp-mass hypothesis. The exact quantitative Herbst assumption is in the cited scattering article rather than in this short report passage; no invented exponent or spectral bound is substituted here.

\paragraph{Catalogue mathematical target.}
\paragraph{Catalogue wording retained for traceability.}
Can Herbst's technical regularity assumption in the extension of Haag–Ruelle scattering theory to theories with massless particles be replaced by a physically motivated criterion?

\paragraph{Saved literature context.} The supplied catalogue's assessment dated 2026-08-21 records: The report explicitly identifies replacing Herbst's technical assumption by a physically motivated criterion as an open problem. The assessment is retained as dated provenance; it is not a new comprehensive literature or open-status certification.
\subsection{Short English statement}
Replace the technical spectral regularity assumption in massless-background Haag--Ruelle scattering by a physically meaningful condition.
\subsection{Research attempt}
\paragraph{Model, date and investigated scope.}
OpenAI GPT-6 Astra (configured model \texttt{gpt-6-astra}), ultra reasoning; generated 11 October 2026. Proves an inverse-spectral-moment sufficient bound and constructs a finite-positive-moment counterexample to power-law regularity.

\paragraph{Sources and literature check.}
The complete supplied definition, including source qualifications and the dated literature assessment, was reviewed. The calculation below is an elementary restricted attempt motivated by [S2]; it does not claim novelty or a fresh literature-status determination.

\paragraph{Context and investigated scope.} The complete supplied definition, including source qualifications and the dated literature assessment, was reviewed. The calculation below is an elementary restricted attempt motivated by [S2]; it does not claim novelty or a fresh literature-status determination.

\paragraph{An abstract spectral test.} Let $K\geq0$ be a self-adjoint operator measuring spectral distance from a chosen sharp shell, let $P$ project onto its kernel, and let $\xi=(I-P)A\Omega$. Write $\mu(B)=\langle\xi,1_B(K)\xi\rangle$. This is an auxiliary abstraction, not a claimed transcription of Herbst's exact condition. If for some $s>0$
\[
 \int_{(0,1)}r^{-2s}\,d\mu(r)<\infty,
\]
then $\|1_{(0,\delta)}(K)\xi\|\leq C\delta^s$ for $0<\delta<1$: on that interval $1\leq\delta^{2s}r^{-2s}$, and integration proves the assertion.

\paragraph{Counterexample to a naive energy criterion.} Finite positive energy moments alone imply no such power law. On $(0,e^{-1})$ put $d\mu(r)=dr/[r(\log(1/r))^2]$. It is a probability measure and
\[
 \mu((0,\delta))=1/\log(1/\delta),\qquad 0<\delta<e^{-1}.
\]
All positive moments are finite because $r$ is bounded; nevertheless this mass exceeds every fixed multiple of $\delta^{2s}$ for all sufficiently small $\delta$. Realize the example by multiplication by $r$ on $L^2(\mu)$ with $\xi=1$, and adjoin a one-dimensional zero eigenspace to provide a sharp eigenvalue. The added shell does not improve the orthogonal component.

\paragraph{Source-scope limit.} The example is a spectral model, not a local relativistic quantum field theory, and says nothing about charged infraparticles, excluded in the source. It rules out the purely abstract implication ``sharp shell plus finite energy moments implies power-law spectral regularity''. Deriving a stronger near-shell bound from physically meaningful locality or dynamical assumptions, and matching it to the exact scattering estimates of the cited article, are the unresolved steps.

\paragraph{Verification and adversarial audit.}
Checked measure normalization, logarithmic tail, bounded positive moments, added sharp eigenspace, and that the toy operator is not asserted to be an AQFT example. Model self-audit only; no independent human proof checking. Prompts A and B were applied to this canonical record and its qualified source scope. Checked measure normalization, logarithmic tail, bounded positive moments, added sharp eigenspace, and that the toy operator is not asserted to be an AQFT example. Retained unresolved status and zero original-target progress; No physically justified criterion has been shown to imply the exact near-shell estimates required by the source scattering theorem.

Coordinator checked normalization and logarithmic small-energy mass of the explicit spectral measure and the inverse-moment sufficient bound. The abstract operator is not represented as a local field theory. This additional check was model-assisted, not independent human verification.

\paragraph{Tools, process and contributions.}
Python 3 was used for file preparation and structural checks; no numerical experiment or theorem prover was used. Contribution prompt A was followed by prompt B. GPT-6 Astra with ultra reasoning drafted the research and performed its stated self-audit; any additional coordinator contribution is identified in the verification and audit. A GPT-6-based Codex coordinator assembled the manuscript and checked repository integration; its exact variant and reasoning configuration were not exposed. No human mathematical contribution, independent human proof check or formal proof certification is supplied. Complete derivations and the consulted source references are included here; private reasoning transcripts are not published.

\paragraph{Final status and first remaining gap.}
\textbf{UNRESOLVED}. No physically justified criterion has been shown to imply the exact near-shell estimates required by the source scattering theorem.

\section{Completion Estimate}
\noindent\textbf{COMPLETION: 0\%}
The calculation is an auxiliary or restricted attempt and establishes no newly verified portion of the full source target. No physically justified criterion has been shown to imply the exact near-shell estimates required by the source scattering theorem.

\subsection{Sources}
\begin{itemize}
\item[\hypertarget{source:statement:1}{S1}] ULAM.AI and the UnsolvedMath Contributors, UnsolvedMath, catalogue entry. Catalogue curation under CC BY 4.0; definitions and English exposition added locally.
\url{https://huggingface.co/datasets/ulamai/UnsolvedMath}.
\item[\hypertarget{source:statement:2}{S2}] Analysis and Quantum Theory (2005), Oberwolfach Reports, DOI 10.4171/owr/2005/44, pp. 2493--2518; official report source and mathematical context.
\url{https://doi.org/10.4171/owr/2005/44}.
\item[\hypertarget{source:statement:3}{S3}] Official EMS Press PDF of the same Oberwolfach report; PDF pages 7--8, printed pages 2499--2500.
\url{https://ems.press/content/serial-article-files/46017}.
\end{itemize}
\label{end:statement}
\end{document}
